\begin{textsample}{POS Dim 7 – human – Score 4.00 – mg/mg\_inf\_e\_ef\_intro\_1\_p11.txt}  \label{ex:f7_pos_008}
A diversidade \textbf{regional} do Estado de Minas Gerais é resultado de um processo histórico de ocupação do território marcado por diferentes fatores, desde aqueles de ordem socioeconômica até os \textbf{naturais} de clima e vegetação.
Essa diversidade se traduz no que podemos entender como várias “ Minas Gerais ” dentro dos limites do Estado, exigindo, portanto, diferentes formas de \textbf{abordagem} e atuação sobre a realidade mineira.
De fato, a efetividade de qualquer iniciativa parte necessariamente da compreensão da realidade para a qual se propõe.

% matched lemmas: abordagem, natural, regional
\end{textsample}
